\section{Mathematische Modellierung}

Im Folgenden wird die geometrische Problemstellung mathematisch beschrieben.
Ziel ist es, zunächst den Flächeninhalt der zusammengesetzten Figur in
Abhängigkeit von der Schenkellänge \(s\) und dem Winkel \(\alpha\) zu
bestimmen. Anschließend wird der Winkel hergeleitet, bei dem dieser
Flächeninhalt maximal wird.

\subsection{Geometrische Vorüberlegungen}

Das betrachtete Dreieck ist gleichschenklig. Die beiden gleich langen
Schenkel besitzen jeweils die Länge \(s>0\). Der eingeschlossene Winkel
zwischen diesen beiden Schenkeln ist \(\alpha\). Da der Winkel im Bereich
von \(0^\circ\) bis \(180^\circ\) variiert werden kann, betrachten wir im
Folgenden

\begin{equation}
    \alpha \in [0,\pi].
\end{equation}

Für die analytischen Rechnungen wird also das Bogenmaß verwendet. Die
Umrechnung in Grad erfolgt über

\begin{equation}
    \alpha_{\degree} = \alpha \cdot \frac{180^\circ}{\pi}.
\end{equation}

Die Basisseite des Dreiecks besitzt die Länge \(2r\). Da das Dreieck
gleichschenklig ist, halbiert die Höhe auf die Basisseite sowohl die
Basisseite als auch den Winkel \(\alpha\). Dadurch entstehen zwei
kongruente rechtwinklige Dreiecke. In einem dieser Teildreiecke ist \(s\)
die Hypotenuse, \(r\) die halbe Basis und der Winkel an der Spitze beträgt
\(\frac{\alpha}{2}\).

Damit gilt

\begin{equation}
    \sin\left(\frac{\alpha}{2}\right)
    =
    \frac{r}{s}.
\end{equation}

Folglich ergibt sich der Radius des angelagerten Halbkreises zu

\begin{equation}
    r
    =
    s\sin\left(\frac{\alpha}{2}\right).
    \label{eq:radius}
\end{equation}

Analog ergibt sich für die Höhe \(h\) des Dreiecks

\begin{equation}
    h
    =
    s\cos\left(\frac{\alpha}{2}\right).
    \label{eq:hoehe}
\end{equation}

Diese beiden Beziehungen bilden die Grundlage für die Berechnung der
Dreiecks- und Halbkreisfläche.

\subsection{Flächenberechnung}

Die Gesamtfläche setzt sich aus der Dreiecksfläche und der Halbkreisfläche
zusammen. Es gilt also

\begin{equation}
    A(\alpha)
    =
    A_{\mathrm{dreieck}}(\alpha)
    +
    A_{\mathrm{halb}}(\alpha).
    \label{eq:gesamtflaeche_summe}
\end{equation}

Zunächst wird die Fläche des Dreiecks betrachtet. Mit der Grundseite
\(2r\) und der Höhe \(h\) gilt

\begin{equation}
    A_{\mathrm{dreieck}}
    =
    \frac{1}{2}\cdot 2r \cdot h
    =
    r h.
\end{equation}

Durch Einsetzen von \eqref{eq:radius} und \eqref{eq:hoehe} folgt

\begin{equation}
    A_{\mathrm{dreieck}}(\alpha)
    =
    s^2
    \sin\left(\frac{\alpha}{2}\right)
    \cos\left(\frac{\alpha}{2}\right).
\end{equation}

Mithilfe der trigonometrischen Identität ergibt sich

\begin{equation}
    A_{\mathrm{dreieck}}(\alpha)
    =
    \frac{1}{2}s^2\sin(\alpha).
    \label{eq:dreiecksflaeche}
\end{equation}

Der an die Basisseite angelagerte Halbkreis besitzt den Radius \(r\).
Da die Fläche eines Kreises mit Radius \(r\) gleich \(\pi r^2\) ist, gilt
für den Halbkreis

\begin{equation}
    A_{\mathrm{halb}}
    =
    \frac{1}{2}\pi r^2.
\end{equation}

Mit \eqref{eq:radius} folgt

\begin{equation}
    A_{\mathrm{halb}}(\alpha)
    =
    \frac{1}{2}\pi s^2
    \sin^2\left(\frac{\alpha}{2}\right).
    \label{eq:halbkreisflaeche}
\end{equation}

Die Gesamtfläche erhält man nun durch Einsetzen von
\eqref{eq:dreiecksflaeche} und \eqref{eq:halbkreisflaeche} in
\eqref{eq:gesamtflaeche_summe}:

\begin{equation}
    A(\alpha)
    =
    \frac{1}{2}s^2\sin(\alpha)
    +
    \frac{1}{2}\pi s^2
    \sin^2\left(\frac{\alpha}{2}\right).
    \label{eq:gesamtflaeche_halbwinkel}
\end{equation}

Zur weiteren Berechnung vereinfachen wir die Formel. Mit

\begin{equation}
    \sin^2\left(\frac{\alpha}{2}\right)
    =
    \frac{1-\cos(\alpha)}{2}
\end{equation}

ergibt sich

\begin{equation}
    A(\alpha)
    =
    \frac{1}{2}s^2\sin(\alpha)
    +
    \frac{1}{4}\pi s^2
    \left(1-\cos(\alpha)\right).
\end{equation}

Durch Ausklammern von \(\frac{s^2}{4}\) folgt die für die Optimierung
verwendete Darstellung

\begin{equation}
    A(\alpha)
    =
    \frac{s^2}{4}
    \left(
        2\sin(\alpha)
        +
        \pi\left(1-\cos(\alpha)\right)
    \right).
    \label{eq:gesamtflaeche}
\end{equation}

\subsection{Winkelberechnung}

Gesucht ist nun der Winkel \(\alpha\), für den die Gesamtfläche $A(\alpha)$ maximal wird. Da \(s>0\) gilt, ist \(s^2\) eine positive Konstante; der optimale Winkel hängt somit nicht von der Schenkellänge \(s\) ab. Zur Bestimmung des Maximums wird \eqref{eq:gesamtflaeche} nach \(\alpha\)
abgeleitet:

\begin{equation}
    A'(\alpha)
    =
    \frac{s^2}{4}
    \left(
        2\cos(\alpha)
        +
        \pi\sin(\alpha)
    \right).
    \label{eq:erste_ableitung}
\end{equation}

Ein inneres Extremum erfüllt

\begin{equation}
    A'(\alpha)=0.
\end{equation}

Da \(\frac{s^2}{4}>0\) gilt, ist dies äquivalent zu

\begin{equation}
    2\cos(\alpha)
    +
    \pi\sin(\alpha)
    =
    0.
\end{equation}

Umstellen liefert

\begin{equation}
    2\cos(\alpha)
    =
    -\pi\sin(\alpha).
\end{equation}

Für \(\alpha\in(0,\pi)\) gilt \(\sin(\alpha)>0\). Daher kann durch
\(\sin(\alpha)\) dividiert werden. Es folgt

\begin{equation}
    \tan(\alpha)
    =
    -\frac{2}{\pi}.
\end{equation}

Da \(\tan(\alpha)<0\) ist und \(\alpha\in(0,\pi)\) gilt, liegt der gesuchte
Winkel im zweiten Quadranten. Somit ergibt sich

\begin{equation}
    \alpha_{\max}
    =
    \pi
    -
    \arctan\left(\frac{2}{\pi}\right).
    \label{eq:alpha_max}
\end{equation}

Numerisch entspricht $\alpha_{\max}$ damit jeweils im Bogen- und Gradmaß:

\begin{equation}
    \alpha_{\max}
    \approx
    2.574681149\ 
    \approx
    147.52^\circ
\end{equation}

Zur Überprüfung, dass es sich tatsächlich um ein Maximum handelt, wird die
zweite Ableitung betrachtet. Aus \eqref{eq:erste_ableitung} folgt

\begin{equation}
    A''(\alpha)
    =
    \frac{s^2}{4}
    \left(
        -2\sin(\alpha)
        +
        \pi\cos(\alpha)
    \right).
\end{equation}

Für den Winkel \(\alpha_{\max}\) liegt \(\alpha\) im zweiten Quadranten.
Deshalb sind sowohl \(\sin(\alpha_{\max})>0\) als auch \(\cos(\alpha_{\max})<0\) und es gilt:

\begin{equation}
    A''(\alpha_{\max})<0.
\end{equation}

Der gefundene Winkel beschreibt also ein lokales Maximum. Da an den
Randpunkten \(\alpha=0\) und \(\alpha=\pi\) keine größere Fläche entsteht,
ist dieses Maximum zugleich das globale Maximum im betrachteten
Winkelbereich.

\subsection{Maximaler Flächeninhalt}

Zur Berechnung des maximalen Flächeninhalts wird
\eqref{eq:alpha_max} in \eqref{eq:gesamtflaeche} eingesetzt. Setze dazu

\begin{equation}
    \beta
    =
    \arctan\left(\frac{2}{\pi}\right).
\end{equation}

Dann gilt

\begin{equation}
    \alpha_{\max}
    =
    \pi-\beta.
\end{equation}

Aus

\begin{equation}
    \tan(\beta)=\frac{2}{\pi}
\end{equation}

folgen

\begin{equation}
    \sin(\beta)
    =
    \frac{2}{\sqrt{\pi^2+4}}
\end{equation}

und

\begin{equation}
    \cos(\beta)
    =
    \frac{\pi}{\sqrt{\pi^2+4}}.
\end{equation}

Da \(\alpha_{\max}=\pi-\beta\) gilt, erhalten wir

\begin{equation}
    \sin(\alpha_{\max})
    =
    \frac{2}{\sqrt{\pi^2+4}}
\end{equation}

und

\begin{equation}
    \cos(\alpha_{\max})
    =
    -\frac{\pi}{\sqrt{\pi^2+4}}.
\end{equation}

Einsetzen in \eqref{eq:gesamtflaeche} liefert

\begin{equation}
    A_{\max}(s)
    =
    \frac{s^2}{4}
    \left(
        2\cdot \frac{2}{\sqrt{\pi^2+4}}
        +
        \pi
        \left(
            1
            -
            \left(
                -\frac{\pi}{\sqrt{\pi^2+4}}
            \right)
        \right)
    \right).
\end{equation}

Nach Zusammenfassen der Terme ergibt sich als gesuchte Berechnungsvorschrift für den maximalen Flächeninhalt

\begin{equation}
    A_{\max}(s)
    =
    \frac{\pi+\sqrt{\pi^2+4}}{4}s^2.
    \label{eq:a_max}
\end{equation}


